\section{Related work}
\label{sec:related}

\subsection{LLM agents in industrial systems}
\label{sec:related-agents}

A growing body of work connects LLM agents to production equipment. At the
execution level, an agent turn interleaves reasoning with tool invocation
\cite{yao2023react}, and models can be taught to call tools
\cite{schick2023toolformer}; the governing question is therefore not
whether an agent can invoke an operation but what bounds that invocation.
Architectural studies integrate LLMs with digital twins and automation
systems \cite{xia2025architecture} and with asset administration shells and
OPC UA information models \cite{graul2026aas}. Control-facing systems use
LLM agents to operate automation systems \cite{xia2025control}, to explore
resource capabilities in multi-agent manufacturing
\cite{lim2025llmenabled}, to drive skill-based manufacturing control
\cite{koehle2026llmdriven}, and to compose semantic recipes on OPC UA
skills \cite{koehle2025semanticrecipes}. Shop-floor orchestration systems
coordinate multi-agent LLM teams for manufacturing tasks
\cite{zhao2026shopfloor}, and pretrained LLMs have been placed in the loop
of robot-operated serial production lines \cite{waseem2026pretrained}.
Process-industry studies use LLM agents with digital twins for fault
handling \cite{gill2025llmdt} and natural-language machine operation over
OPC UA \cite{hofmann2025nlc}. Reviews of agentic AI for process monitoring
and fault diagnosis \cite{jiang2026agentic} and of LLM-enabled digital
twins \cite{gill2025llmdt} catalog the application space. Closer to control
theory, InstructMPC embeds human--LLM interaction into context-aware MPC
\cite{wu2025instructmpc}, and LLM-based architectures have been proposed
for automated control-code generation in industrial cyber-physical systems
\cite{huang2026made} and closed-loop PLC code generation with verification
\cite{liu2026agents4plc}. This journal has hosted platform papers of
comparable shape, including an open-source compositional digital twin
framework \cite{opentwins2023} and requirements and design patterns for
adaptive digital twins \cite{dtpatterns2023}, as well as a controlled
assessment of an LLM-based digital intelligent assistant in assembly
manufacturing \cite{llmdia2024}.

Across these systems the LLM occupies one of two positions. In
\emph{advisory} positions the agent reads state and recommends; a human or a
conventional controller acts. In \emph{actuating} positions the agent's
output reaches equipment directly \cite{waseem2026pretrained,
xia2025control,koehle2026llmdriven,hofmann2025nlc}. The actuating position is
where governance matters most, and it is the position this paper addresses.

\subsection{Governing machine-originated actions}
\label{sec:related-gov}

Human-factors research frames the delegation of action to automation as a
allocation problem across types and levels of human interaction
\cite{parasuraman2000model}, and warns that human oversight of a
predominantly automatic system degrades through complacency and bias
\cite{parasuramanmanzey2010complacency,bainbridge1983ironies}. These
results motivate designs in which the human decision is exercised against a
concrete, inspectable proposal rather than a summary. In the agent-safety
literature, runtime enforcement systems constrain LLM agent behavior with
user-defined rules \cite{wang2025agentspec}; capability-based architectures
defeat prompt injection by construction \cite{debenedetti2025camel};
benchmarks formalize injection attacks and defenses
\cite{liu2024pinjection,zhan2024injecagent} and dynamic environments
evaluate attacks and defenses against tool-using agents
\cite{debenedetti2024agentdojo}, demonstrating that agents that read
plant-facing content are exposed to manipulated instructions. Industrial
security standards contribute the zone-and-conduit model
\cite{iec62443-3-2} and system security requirements
\cite{iec62443-3-3}, and safety instrumented systems define the
plant-side final barrier \cite{iec61511}; both assume the information
system's requests arrive through designed channels, which is exactly the
assumption an unconstrained agent violates. Architectural work on
digital-twin-based automation security \cite{gehrmann2020dtiacs}, secure
fog layers for the industrial IoT \cite{sengupta2021securefog}, and
functional architectures for special-purpose industrial networks
\cite{mahmood2023functional} shows that separating concerns across layers
is an established design idiom for industrial trust boundaries.

On the industrial side, Lin et al.~\cite{lin2026agentic} integrate policy-governed
agentic AI with multi-source sensor--MES--ERP data for closed-loop
manufacturing control, and Schall~\cite{schall2026safe} validate LLM proposals
against P\&ID-derived structure before they reach process control.
Biskupovic et al.~\cite{biskupovic2025supervision} realize agent-based supervision for
service-oriented industrial cyber-physical systems. These works establish
that plant writes need governance. What remains under-specified is the
audit dimension: the ability to answer, after the fact, \emph{which layer}
of the path admitted or rejected a specific write, against \emph{which
configuration revision}, and with \emph{what device-side receipt}. Existing
descriptions typically report governance as a property of the overall
pipeline; they do not separate the layers, nor grade the evidence that a
write actually took effect.

Two honest boundaries on this positioning should be stated. First,
\emph{the individual mechanisms have industrial precedents}: limit
checking against engineered ranges, post-write verification, and
approval workflows with audit trails are established practice in
process control and regulated industries (cf. electronic-record audit
trail requirements such as 21 CFR Part 11), and run-to-run control
already bounded recipe execution decades ago \cite{moyne2000r2r}. The
contribution claimed here is not any single mechanism but their
configuration-driven composition into one agent-facing path in which
every layer's verdict is separately recorded, the deciding layer of a
rejection is identifiable, receipts are graded by device-side evidence,
and no bypass exists from the agent surface. Second, \emph{a conservative
alternative exists}: keeping the LLM strictly read-only and leaving
writes to conventional DCS limits plus historian audit covers much of
the same ground when actuating agents are not wanted. The platform's
value proposition is for deployments that \emph{do} want supervised
agent-initiated writes, where it makes such writes governed by
construction rather than by convention. Table~\ref{tab:position}
positions representative systems along this dimension.

\begin{table}[!t]
\centering
\caption{Representative LLM-agent industrial systems and their treatment of
plant-facing writes (characterizations at citation level, from the cited
works' stated scope).}
\label{tab:position}
\small
\begin{spacing}{1.0}
\begin{tabularx}{\columnwidth}{@{}p{0.24\columnwidth}p{0.30\columnwidth}X@{}}
\toprule
Study & LLM role & Plant-facing write treatment \\
\midrule
Zhao et al.~\cite{zhao2026shopfloor} & Shop-floor multi-agent orchestration & Task-level coordination; not centered on setpoint writes \\
Xia et al.~\cite{xia2025control} & Operates automation systems & Direct system operation; engineering limits not the focus \\
K{\"o}hle et al.~\cite{koehle2026llmdriven} & Skill-based control & Skill invocation; skills encapsulate equipment access \\
Lin et al.~\cite{lin2026agentic} & Policy-governed closed-loop control & Policies govern agent actions in the control loop \\
Schall~\cite{schall2026safe} & Process-control integration & P\&ID-grounded validation of proposals \\
Wu et al.~\cite{wu2025instructmpc} & Context provider for MPC & Human--LLM interaction around a conventional controller \\
Hofmann et al.~\cite{hofmann2025nlc} & Natural-language operation & NL commands mapped to OPC UA operations \\
\midrule
This work & Supervisory agent teams & Layer-attributable interlocks, bound approvals, receipt grading, audited ledger \\
\bottomrule
\end{tabularx}
\end{spacing}
\end{table}
